\begin{theorem}[MellinConvolutionTransform]
\uses{MellinConvolution}
\label{MellinConvolutionTransform}
\lean{MellinConvolutionTransform}
% at /home/runner/work/trough/trough/library/PrimeNumberTheoremAnd/MellinCalculus.lean:357.0-432.26
\leanok
Let $f$ and $g$ be functions from $\mathbb{R}_{>0}$ to $\mathbb{C}$ such that
  \begin{equation}
    (x,y)\mapsto f(y)\frac{g(x/y)}yx^{s-1}
    \label{eq:assm_integrable_Mconv}
  \end{equation}
  is absolutely integrable on $[0,\infty)^2$.
  Then
  $$\mathcal{M}(f\ast g)(s) = \mathcal{M}(f)(s)\mathcal{M}(g)(s).$$
\end{theorem}
\begin{proof}
\leanok
By Definitions \ref{MellinTransform} and \ref{MellinConvolution}
  $$
    \mathcal M(f\ast g)(s)=
    \int_0^\infty \int_0^\infty f(y)g(x/y)x^{s-1}\frac{dy}ydx
  $$
  By (\ref{eq:assm_integrable_Mconv}) and Fubini's theorem,
  $$
    \mathcal M(f\ast g)(s)=
    \int_0^\infty \int_0^\infty f(y)g(x/y)x^{s-1}dx\frac{dy}y
  $$
  in which we change variables from $x$ to $z=x/y$:
  $$
    \mathcal M(f\ast g)(s)=
    \int_0^\infty \int_0^\infty f(y)g(z)y^{s-1}z^{s-1}dzdy
  $$
  which, by Definition \ref{MellinTransform}, is
  $$
    \mathcal M(f\ast g)(s)=
    \mathcal M(f)(s)\mathcal M(g)(s)
    .
  $$
\end{proof}
